% ------------------------------------------------------------------
%  Capitolo 1 — capitolo di esempio: mostra tutti gli elementi
%  della classe. Sostituisci il testo con il tuo.
% ------------------------------------------------------------------
\chapter{Non per la scuola, ma per la vita}
\label{cap:vita}

\epigrafe{Non vitae sed scholae discimus: impariamo non per la vita,
  ma per la scuola.}{Seneca}

\iniziale{S}{i dice spesso} \enquote{non scholae sed vitae discimus}:
non impariamo per la scuola, ma per la vita. Pochi sanno che la frase
originale di Seneca dice l'esatto contrario ed è un rimprovero: il
filosofo lamentava che i suoi contemporanei studiassero per le
dispute delle scuole, non per vivere meglio.\footnote{Seneca,
\emph{Epistulae morales ad Lucilium}, 106, 12. La versione rovesciata,
diventata un motto scolastico, è di molti secoli successiva.} Il
rimprovero è ancora attuale.

Questo capitolo di esempio serve a mostrare gli elementi tipografici
della classe: paragrafi, note, citazioni, riquadri, figure,
riferimenti bibliografici e la sintesi di fine capitolo.

\section{La curva dell'oblio}

Alla fine dell'Ottocento Hermann Ebbinghaus misurò su sé stesso la
velocità con cui dimenticava liste di sillabe prive di senso
\parencite{ebbinghaus1885}. Il risultato è la celebre \emph{curva
dell'oblio}: la maggior parte di ciò che impariamo si perde nelle prime
ore e nei primi giorni, poi la perdita rallenta
(figura~\ref{fig:oblio}).

\begin{figure}[tbp]
  \centering
  \begin{tikzpicture}[x=0.95cm, y=2.6cm, font=\footnotesize]
    \draw[filetto!70, line width=0.3pt] (0,0) -- (6.4,0);
    \draw[filetto!70, line width=0.3pt] (0,0) -- (0,1.08);
    \foreach \x/\l in {0/0, 1/1, 2/2, 3/3, 4/4, 5/5, 6/6}
      \draw[filetto!70] (\x,0) -- (\x,-0.03) node[below] {\l};
    \foreach \y/\l in {0/0, 0.5/50, 1/100}
      \draw[filetto!70] (0,\y) -- (-0.08,\y) node[left] {\l};
    \node[below] at (3.2,-0.16) {\itshape giorni dall'apprendimento};
    \node[rotate=90, above] at (-0.55,0.5) {\itshape \% ricordato};
    % curva senza ripasso
    \draw[inchiostro, line width=0.6pt, domain=0:6.2, samples=80]
      plot (\x, {0.25 + 0.75*exp(-1.6*\x)});
    % curva con ripassi (giorni 1, 3)
    \draw[rubrica, line width=0.6pt, domain=0:1, samples=30]
      plot (\x, {0.25 + 0.75*exp(-1.6*\x)});
    \draw[rubrica, line width=0.6pt, dashed] (1,0.40) -- (1,1);
    \draw[rubrica, line width=0.6pt, domain=1:3, samples=40]
      plot (\x, {0.5 + 0.5*exp(-0.8*(\x-1))});
    \draw[rubrica, line width=0.6pt, dashed] (3,0.60) -- (3,1);
    \draw[rubrica, line width=0.6pt, domain=3:6.2, samples=40]
      plot (\x, {0.7 + 0.3*exp(-0.35*(\x-3))});
    \node[rubrica, anchor=west] at (4.1,0.93) {\itshape con ripasso};
    \node[anchor=west] at (3.9,0.18) {\itshape senza ripasso};
  \end{tikzpicture}
  \caption{La curva dell'oblio (andamento qualitativo) e l'effetto
    di due ripassi distanziati nel tempo.}
  \label{fig:oblio}
\end{figure}

La buona notizia è che ogni ripasso ben distanziato rende la curva
più piatta: dimentichiamo di meno, e più lentamente. È il cosiddetto
\emph{effetto di spaziatura}, uno dei risultati più solidi della
psicologia sperimentale \parencite{cepeda2006}.

\begin{daricordare}
Ciò che non viene richiamato alla memoria si perde in fretta. Ciò che
viene richiamato più volte, a intervalli crescenti, resta.
\end{daricordare}

\subsection{Rileggere non basta}

Molti studenti ripassano rileggendo e sottolineando. Sono le strategie
più diffuse e, secondo un'ampia rassegna della letteratura, fra le meno
efficaci \parencite{dunlosky2013}. Rileggere dà una sensazione di
familiarità che viene scambiata per conoscenza.

\begin{quote}
Qui trova posto una citazione lunga, composta in corpo minore e
rientrata su entrambi i lati, così da staccarsi dal testo senza
bisogno di virgolette.
\end{quote}

\asterismo

\section{Interrogarsi per ricordare}

Un esperimento divenuto classico ha confrontato due gruppi di studenti:
il primo rileggeva un testo più volte, il secondo lo leggeva una volta
e poi provava a richiamarlo a memoria. A distanza di una settimana il
secondo gruppo ricordava molto di più \parencite{roediger2006}. È
l'\emph{effetto di verifica} (in inglese \emph{testing effect}),
discusso a lungo in \textcite{brown2014}.

\begin{inpratica}
\begin{enumerate}
  \item Chiudi il libro dopo ogni paragrafo.
  \item Scrivi, senza guardare, ciò che ricordi.
  \item Riapri il libro e confronta: ciò che mancava è ciò da studiare.
\end{enumerate}
\end{inpratica}

\subsubsection{Una tabella} Anche le tabelle seguono lo stile
classico, con filetti sottili e senza linee verticali.

\begin{table}[tbp]
  \centering\small
  \begin{tabular}{@{}lc@{}}
    \toprule
    \textit{Strategia} & \textit{Utilità} \\
    \midrule
    Verifica (richiamo attivo) & alta \\
    Ripasso distanziato       & alta \\
    Rilettura                 & bassa \\
    Sottolineatura            & bassa \\
    \bottomrule
  \end{tabular}
  \caption{Alcune strategie di studio e la loro utilità secondo
    \textcite{dunlosky2013}.}
  \label{tab:strategie}
\end{table}

\begin{domande}
  \item Che cosa mostra la curva dell'oblio?
  \item Perché la rilettura dà un'illusione di conoscenza?
  \item Come si può applicare l'effetto di verifica a una materia
        che studi adesso?
\end{domande}

\begin{sintesi}
  \item Si dimentica in fretta ciò che non viene richiamato.
  \item Il ripasso distanziato appiattisce la curva dell'oblio.
  \item Interrogarsi vale più che rileggere.
\end{sintesi}

\finecapitolo
